%=====================================================================
\chapter{Quasi-Experiments and Repository Mining}\label{ch:quasiexp}
%=====================================================================
\locator{3}{Part~3 --- Research Designs for Computing}

\begin{outcomes}
By the end of this chapter you will be able to:
\begin{tight}
  \item \textbf{Explain} the selection problem and why it, not sample size, is
    the central difficulty in observational research.
  \item \textbf{Choose} an identification strategy --- difference-in-differences,
    interrupted time series, regression discontinuity or matching --- and state
    the assumption it rests on.
  \item \textbf{Design} a repository mining study and enumerate its data
    quality threats.
  \item \textbf{Judge} when an observational claim may legitimately be phrased
    causally.
\end{tight}
\end{outcomes}

\begin{prereq}
\chref{experiments}, especially what randomisation buys. This chapter is about
what to do when you cannot have it. \chref{questions} for confounders.
\end{prereq}

\begin{keyterms}
quasi-experiment $\bullet$ selection bias $\bullet$ identification strategy
$\bullet$ counterfactual $\bullet$ difference-in-differences $\bullet$ parallel
trends $\bullet$ interrupted time series $\bullet$ regression discontinuity
$\bullet$ propensity score matching $\bullet$ instrumental variable $\bullet$
natural experiment $\bullet$ mining software repositories $\bullet$
survivorship bias $\bullet$ tangled commit
\end{keyterms}

\section{The Problem Is Selection, Not Size}

Most interesting questions in software engineering cannot be answered by
experiment. You cannot randomly assign companies to adopt microservices, or
developers to have ten years of experience. What you can do is observe what
happened where the choice was made by someone else --- and the difficulty is
that the choice was not made at random.

\begin{definitionbox}[The selection problem]
When units choose their own treatment, the treated and untreated groups differ
in every respect that influenced the choice --- including respects you cannot
observe.

\medskip
A comparison of outcomes then confounds the effect of the treatment with the
effect of \emph{being the kind of unit that selects into it}. More data does
not help: a million observations estimate the confounded quantity with great
precision.
\end{definitionbox}

\begin{conceptbox}[What an identification strategy is]
An \term{identification strategy} is an argument --- not a statistical
technique --- that the comparison you are making approximates the
counterfactual: what would have happened to the treated units had they not been
treated.

\medskip
Every strategy below rests on an assumption that cannot be tested directly.
Your job in the thesis is to state that assumption plainly, give whatever
indirect evidence for it you can, and say what would happen to your conclusion
if it failed. A quasi-experimental study whose assumption is stated and
defended is respectable research. One that runs a regression and says
``controls were included'' is not.
\end{conceptbox}

\section{Four Strategies}

\begin{center}
\small
\begin{longtable}{@{}L{2.7cm}L{4.6cm}L{6.4cm}@{}}
\toprule
\textbf{Strategy} & \textbf{How it identifies} & \textbf{The assumption it rests on} \\
\midrule
\endfirsthead
\toprule
\textbf{Strategy} & \textbf{How it identifies} & \textbf{The assumption it rests on} \\
\midrule
\endhead
\bottomrule
\endfoot
Difference-in\-differences
  & Compares the \emph{change} in treated units with the \emph{change} in
    untreated units over the same period
  & \term{Parallel trends}: absent the treatment, both groups would have
    changed by the same amount. Partly checkable by inspecting pre-treatment
    trends \\
Interrupted time series
  & Models the series before an intervention and tests for a change in level
    or slope at the intervention point
  & No other change occurred at the same moment, and the pre-trend would have
    continued \\
Regression discontinuity
  & Compares units just above and just below a threshold that determines
    treatment
  & Units cannot precisely manipulate their position relative to the
    threshold; nothing else changes at it \\
Propensity score matching
  & Matches treated to untreated units with similar probability of being
    treated, given observed covariates
  & \term{No unobserved confounding} --- everything driving selection was
    measured. Strong, untestable, and usually false \\
\end{longtable}
\end{center}

\begin{alertbox}[Matching is the weakest of the four, and the most used]
Propensity score matching is popular because it is easy to run and produces
balanced-looking tables. But it only balances what you \emph{measured}. If
teams that adopt a practice differ in engineering culture, and you did not
measure engineering culture, matching on repository age and star count changes
nothing about that.

\medskip
Difference-in-differences is usually available in software repository data ---
you almost always have a time series --- and it is far stronger, because each
unit acts as its own control and any time-invariant unmeasured difference
cancels. When you have a before and an after, prefer it.
\end{alertbox}

\begin{worked}[Repairing the static analysis study]
\chref{appraisal} appraised a constructed paper claiming that adopting static
analysis reduces post-release defects, comparing 100 adopting projects with 100
non-adopting ones. The design supports association only. Here is the repair.

\medskip
\textbf{The counterfactual we need.} What would have happened to the adopting
projects' defect rates had they not adopted?

\medskip
\textbf{Difference-in-differences.} For each adopting project, identify the
adoption date. Measure the outcome for 12 months before and 12 months after.
Do the same for non-adopters over a matched calendar window. The estimate is
\[
\hat\delta \;=\; (\bar Y^{\text{post}}_{\text{adopt}} - \bar Y^{\text{pre}}_{\text{adopt}})
              - (\bar Y^{\text{post}}_{\text{control}} - \bar Y^{\text{pre}}_{\text{control}}).
\]

\medskip
\textbf{What this fixes.} Any \emph{time-invariant} difference between adopters
and non-adopters --- better developers, stricter culture, a more testable
architecture --- differences out. That is precisely the confounding the
original design could not address.

\medskip
\textbf{What it does not fix.} Anything that changed \emph{at the same time} as
adoption. If teams adopt static analysis as part of a broader quality push that
also introduced code review and CI, difference-in-differences attributes the
whole package to static analysis. This must be stated as the principal threat,
and can be partly probed by checking whether other practices appear in the same
commits.

\medskip
\textbf{Checking parallel trends.} Plot both groups' monthly defect rates for
24 months before adoption. If the curves are roughly parallel, the assumption
is plausible. If adopters were already improving faster, it is not, and the
estimate is biased --- report the plot either way.

\medskip
\textbf{The construct problem remains.} Issues labelled \texttt{bug} still
measure labelling as much as defects (\chref{measurement}). Difference-in-differences
helps here too, but only if labelling practice did not change at adoption ---
which, if the quality push was real, it probably did. Consider a second outcome
less sensitive to labelling, such as post-release reverts.

\medskip
All figures illustrative; this is a design, not a result.
\end{worked}

\section{Mining Software Repositories}

Software development leaves an unusually rich trace: version control, issue
trackers, code review, CI logs, package registries. Mining it is cheap, scales
to thousands of projects, and is entirely unobtrusive. It is also, as data, far
messier than it appears.

\begin{paperbox}[What repository data will and will not tell you]
Eirini Kalliamvakou, Georgios Gousios, Kelly Blincoe, Leif Singer, Daniel M.
German and Daniela Damian, \emph{The promises and perils of mining GitHub},
MSR 2014~\cite{kalliamvakou2014promises}.

\medskip
\textbf{Why it matters.} It is the standard reference for the data quality
threats in this design, and it is specific: not ``be careful'' but a catalogue
of named perils with evidence for each.

\medskip
\textbf{What to extract.} The list of perils, and the corresponding mitigation
for each one you cannot avoid. Cite it in your threats section and demonstrate
you handled the ones that apply.
\end{paperbox}

\begin{center}
\small
\begin{longtable}{@{}L{3.0cm}L{5.2cm}L{6.0cm}@{}}
\toprule
\textbf{Threat} & \textbf{What goes wrong} & \textbf{What to do} \\
\midrule
\endfirsthead
\toprule
\textbf{Threat} & \textbf{What goes wrong} & \textbf{What to do} \\
\midrule
\endhead
\bottomrule
\endfoot
Not a project
  & Many repositories are course assignments, dotfiles, forks or personal
    experiments
  & Filter on activity, contributor count, issue use and commit history length;
    report the filter \\
Survivorship
  & Abandoned and deleted projects vanish. You study the ones that survived
  & Be explicit that inference is to surviving projects; consider archived
    snapshots \\
Tangled commits
  & A single commit mixes a bug fix, a refactor and a formatting change
  & Do not treat one commit as one change; consider untangling heuristics and
    report their error rate \\
Bot activity
  & Dependabot and CI bots generate large volumes of commits and issues
  & Detect and exclude bot accounts; report how many were removed \\
Missing links
  & Only some commits reference the issue they fix, and not at random
  & Acknowledge that linked-fix datasets are biased towards disciplined
    projects \\
Label semantics
  & \texttt{bug} means different things across projects
  & Validate a sample manually; report agreement \\
Activity is not effort
  & Commit counts measure committing, not work
  & Choose constructs closer to the question (\chref{measurement}) \\
\end{longtable}
\end{center}

\begin{pitfall}[Generalising from GitHub to software engineering]
Open-source projects on one platform are not a random sample of software
development. They differ in ownership, incentives, review practice, team
stability and above all in what is \emph{visible}: much industrial engineering
leaves no public trace.

\medskip
This is a limit on the population your inference reaches
(\chref{validity}), and the fix is not more repositories --- it is stating the
population honestly. ``In active open-source projects on GitHub between 2019
and 2024'' is a defensible scope. ``In software projects'' is not.
\end{pitfall}

\section{When May You Say ``Causes''?}

\begin{center}
\small
\begin{tabular}{@{}L{4.0cm}L{10.3cm}@{}}
\toprule
\textbf{Design} & \textbf{Language it licenses} \\
\midrule
Randomised experiment
  & ``causes'', ``increases'', ``reduces'' --- within the studied population \\
Strong quasi-experiment with a defended assumption
  & ``is associated with, and the difference-in-differences estimate is
    consistent with a causal effect of X on Y under parallel trends'' \\
Matching on observed covariates
  & ``is associated with, after adjusting for observed characteristics'' ---
    and say what is unobserved \\
Plain correlation
  & ``is associated with''. Nothing more, in the title as well as the text \\
\bottomrule
\end{tabular}
\end{center}

\begin{conceptbox}[The title is part of the claim]
A paper whose body carefully says ``associated with'' and whose title says
``improves'' has made a causal claim. Titles are what get read, cited and
repeated. \reg{PhD Regs 2024}{\S14.3(C)} requires that ``the obtained results
should support the study objectives'' --- a title outrunning the design is
exactly the mismatch that clause is aimed at.
\end{conceptbox}

\begin{traditions}{observational study}
\tradEMP{The natural home. The entire methodological effort goes into
constructing a credible counterfactual, and the reporting standard is that a
reader can judge the identifying assumption for themselves.}
\tradDSR{Repository data characterises the problem space before an artefact is
designed, and can supply the naturalistic evaluation baseline afterwards.}
\tradQUAL{Trace data shows what happened but never why. Pairing mining with
interviews of the developers involved is one of the strongest mixed designs
available in our field (\chref{mixed}).}
\tradFORM{Treats the artefacts as formal objects --- dependency graphs, call
graphs, change histories --- where properties can be computed exactly rather
than estimated.}
\end{traditions}

\begin{lab}[Turn a correlation into a design]
Find a published claim in your area that is correlational but written
causally --- they are not hard to find.

\begin{tightnum}
  \item State the counterfactual the claim requires.
  \item Choose an identification strategy that could plausibly recover it from
    available data.
  \item Write the identifying assumption in one sentence, and say what evidence
    you could offer for it.
  \item Name the confounder that would survive your strategy, and say which
    direction it would bias the estimate.
  \item Rewrite the paper's title so that it matches what the original design
    supports.
\end{tightnum}
\end{lab}

\begin{checkpoint}
\begin{tightnum}
  \item Why does increasing the sample size not solve selection bias?
  \item What does difference-in-differences remove that matching does not, and
    what does it fail to remove?
  \item Give two reasons a dataset of commits linked to issues is a biased
    sample of bug fixes.
  \item Under what circumstances, if any, would you use propensity score
    matching in preference to difference-in-differences?
\end{tightnum}
\end{checkpoint}

\begin{chaptersummary}
\begin{tight}
  \item The obstacle in observational research is selection, not sample size.
    Precision does not fix bias.
  \item An identification strategy is an argument about a counterfactual, and
    every one rests on an untestable assumption. State it, defend it, and say
    what happens if it fails.
  \item Prefer difference-in-differences where a before and after exist: it
    removes all time-invariant unmeasured differences. Matching removes only
    what you measured.
  \item Repository data is cheap and scales, and carries specific known threats
    --- non-projects, survivorship, tangled commits, bots, missing links, label
    semantics. Handle the ones that apply and report how.
  \item GitHub is not a sample of software engineering. State the population
    your inference actually reaches.
  \item Match the causal strength of your language --- in the title as well as
    the text --- to the design.
\end{tight}
\end{chaptersummary}

\begin{reviewq}
  \item Parallel trends cannot be verified for the post-treatment period, only
    made plausible for the pre-treatment period. Does this make
    difference-in-differences unscientific? Answer with reference to
    \chref{philosophy}.
  \item Mining studies routinely analyse thousands of projects and report very
    small p-values. Explain why this is not reassuring, and what should be
    reported instead.
  \item Design a natural experiment in software engineering: identify a real
    event that assigned something close to randomly, and say what it could
    identify.
  \item Industrial software leaves little public trace, so the mining
    literature describes open source. What does this systematically distort in
    the field's collective picture of software engineering?
\end{reviewq}
